\section*{4 Rate amplification for dual-induced SLR}
In this section we introduce the notion of smooth locally recoverable sets (SLR) which under non-adaptive and linearity assumptions is shown to be equivalent to LCC. We consider a certain class of SLR, to which we call dual-induced SLR. These are SLR that are obtained by the dual of certain structured sets. The structure of these dual-SLR sets allows us to devise a rate amplification procedures for them. Informally, dual-SLR are sets of tuples (or linear spaces of vectors if the alphabet over which we are working is a field) in which every given entry of a tuple in the set can be recovered using only few queries and in a ``smooth'' manner, which is to say that the distribution of every query has high min-entropy.
\setcounter{nttheorem}{20}
\begin{ntdefinition}[Author Definition21: Smooth locally recoverable sets (SLR)]
Let $\Sigma,P$ be arbitrary non-empty sets. We say that $C\subseteq\Sigma^P$ is $(q,\tau,\varepsilon)$-smooth locally recoverable (SLR for short) if there exists a randomized algorithm $\operatorname{Rec}$, called a recovering procedure, that is given as input $p\in P$ and an oracle access to $c\in C$. The recovering procedure outputs either an element of $\Sigma$ or a symbol $\bot$ which is assumed not to be in $\Sigma$. The algorithm $\operatorname{Rec}$ has the following properties:
\begin{itemize}
\item For every $(c,p)\in C\times P$, $\operatorname{Rec}^c(p)$ makes at most $q$ queries to $c$.
\item For every $c\in C$ and $p,r\in P$ it holds that $\Pr[\operatorname{Rec}^c(p)\text{ queries }c_r]\leq\tau$.
\item For every $(c,p)\in C\times P$, the random variable $\operatorname{Rec}^c(p)\in\{c_p,\bot\}$, and $\Pr[\operatorname{Rec}^c(p)=\bot]\leq\varepsilon$.
\end{itemize}
We assume that for every $p\in P$ whether $\operatorname{Rec}^c(p)=\bot$ is independent of $c$, and that it is never the case that $\operatorname{Rec}^c(p)$ queries $c_p$. When $\Sigma$ is a field and $C$ is a linear subspace of $\Sigma^P$, we say that $C$ is linear. In this case, the rate of $C$ is defined as $\dim(C)/|P|$. We will mostly consider non-adaptive SLR. These are SLR in which the joint distribution of queries is independent of $c$.
\end{ntdefinition}
We remark that the notion of SLR is very similar to the notion of smooth-codes of~\cite{ref23} for LDC. We now have the following easy claim showing that SLR yield LCC and, assuming linearity, LDC.
\begin{ntclaim}[Author Claim22]
Let $C\subseteq\Sigma^P$ be a $(q,\tau,\varepsilon)$-SLR. Then, for every $\varepsilon'>0$, $C$ is a $(q,\delta,\varepsilon+\varepsilon')$-LCC with $\delta=\varepsilon'/(q\tau|P|)$. As a consequence, if $C$ is also linear then $C$ is a $(q,\delta,\varepsilon+\varepsilon')$-LDC.
\end{ntclaim}
\begin{proof}
To show that $C$ is an LCC, we devise a local corrector for $C$. Given an oracle access to $c\in\Sigma^P$, and $p\in P$ as input, the local corrector computes $z=\operatorname{Rec}^c(p)$. If $z=\bot$ then the local corrector returns some arbitrary element of $\Sigma$, and otherwise return $z$. To analyze this local corrector, let $c'\in\Sigma^P$ be such that $\operatorname{dist}(c,c')\leq\delta|P|$. Denote $B=\{p\in P\mid c_p\ne c'_p\}$. Note that conditioned on $\operatorname{Rec}^c(p)\ne\bot$, the local corrector returns $c_p$ successfully if all $q$ queries do not fall into $B$. The probability that any given query falls into $B$ is bounded above by $\tau|B|$ and so, by the union bound, the probability that some query falls into $B$ is bounded above by $\tau|B|q\leq\varepsilon'$. This proves that $C$ is a $(q,\delta,\varepsilon+\varepsilon')$-LCC. Note that linear LCC are systematic and so every linear LCC induces an LDC.
\end{proof}
In fact, for linear LCCs, the other direction also holds, meaning that such an LCC is an SLR, as we have in the following claim.
\begin{ntclaim}[Author Claim23]
Let $C\subseteq\mathbb F^P$ be a linear non-adaptive $(q,\delta,\varepsilon)$-LCC where $1-\varepsilon>1/|\mathbb F|$. Then, $C$ is a $(q,\tau,\varepsilon')$-SLR with $\tau=q/(\delta|P|)$ and $\varepsilon'=0$.
\end{ntclaim}
Before we prove the claim, we need to state the following easy to verify fact.
\begin{nttheorem}[Author Fact24]
Let $L\subseteq\mathbb F^P$ be a linear subspace, let $p\in P$, let $Q\subseteq P$ and let $x\in\mathbb F^{|Q|}$. Then, one of the following cases must hold.
\begin{enumerate}
\item There is at most one $\alpha\in\mathbb F$ for which there exists some $v\in L$ satisfying $v(Q)=x$\footnote{For a set $A=\{a_1,\ldots,a_{|A|}\}$, $v(A)$ denotes the sequence $(v(a_1),\ldots,v(a_{|A|}))$.} and $v(p)=\alpha$;
\item For every $\alpha\in\mathbb F$ there is an equal number of $v\in L$ for which $v(Q)=x$ and $v(p)=\alpha$.
\end{enumerate}
In particular, either no function (even randomized) of $v(Q)$ can predict $v(p)$ with probability more than $1/|\mathbb F|$, when $v\in L$ is randomly chosen uniformly, or $v(Q)$ always determines $v(p)$.
\end{nttheorem}
With that, we now prove Claim23\footnote{This proof is inspired by the proof of~\cite{ref23} of their Theorem1 and by a proof in~\cite{ref41} for a different claim.}.
\begin{proof}[Author proof for Claim23]
To show that $C$ is an SLR, we devise a recovering procedure $\operatorname{Rec}$ for it, based on the local corrector promised by it being an LCC. Let $D$\footnote{We make the slight assumption that $D^c(p)$ never directly queries $c(p)$. If however it does, then similarly $C$ can be shown to be a $(q,\tau,\varepsilon')$-SLR for $\tau=1/((\delta|P|/q)-1)$.} be such a local corrector. For every point $p\in P$, we construct a sequence of disjoint sets $Q_p^1,\ldots,Q_p^{m_p}\subseteq P$, where for every $i$, $c(Q_p^i)$ determines $c(p)$ while satisfying $|Q_p^i|\leq q$, and $m_p\geq\delta|P|/q$. On $p\in P$ and oracle access to $c\in C$, the procedure $\operatorname{Rec}^c(p)$ acts by uniformly choosing $i\in[m_p]$, querying $c(Q_p^i)$, and using it to deduce and output $c(p)$. The correctness of the result of $\operatorname{Rec}$ is immediate (since $\operatorname{Rec}$ always succeeds, $\varepsilon'=0$), and indeed the number of queries is no more than $q$. Since the sets are disjoint, the probability that a point is queried is no more than $\tau=q/(\delta|P|)$. It only remains to show how the assumed sets can be constructed, to conclude that $C$ is a $(q,\tau,\varepsilon')$-SLR, which we now turn to do.
For every $p\in P$, $Q_p^1,\ldots,Q_p^{m_p}$ are constructed as follows. Set $Q_p^0=\varnothing$. For $i=1,2,\ldots$, set $S_i=Q_p^0\cup\cdots\cup Q_p^{i-1}$. If $|S_i|>\delta|P|$, halt and set $m_p=i-1$. Otherwise, it holds that for every $c\in C$, for every modification of the coordinates in $S_i$ to some erroneous values, the decoder $D$ correctly outputs $c(p)$ with probability at least $1-\varepsilon$. An equivalent description of this case is the following: for every $c\in C$ and $z:S_i\to\mathbb F$, define $c_z\in\mathbb F^P$ such that for $r\notin S_i$, $c_z(r)=c(r)$, and for $r\in S_i$, $c_z(r)=z(r)$; the decoder $D$ chooses a set of queries $Q\subseteq P$, $|Q|\leq q$, according to a distribution and applies a function $f_Q$ on $c_z(Q)$; with probability at least $1-\varepsilon$, $f_Q(c_z(Q))=c(p)$. Since $Q$ is chosen in a manner independent of $c$ and $z$, one can verify that this implies that there exists some fixed $Q$ for which when $c\in C$ and $z:\mathbb F\to S_i$ are chosen randomly in a uniform manner, with probability at least $1-\varepsilon$ (this time over the choice of $c$ and $z$), $f_Q(c_z(Q))=c(p)$. Therefore, we can define another function $f'_Q$ that only gets $c(Q\setminus S_i)$, chooses $z$ at random, and outputs $f_Q(c_z(Q))$. If $c\in C$ is chosen uniformly at random, $f'_Q(Q\setminus S_i)=c(p)$ with probability at least $1-\varepsilon>1/|\mathbb F|$. By Fact24, this implies that $c(Q\setminus S_i)$ determines $c(p)$, for every $c\in C$. We therefore set $Q_p^i=Q\setminus S_i$, and proceed to next $i$. As this process only halts when $|S_i|>\delta|P|$, and for every $i$ $|S_i|\leq q(i-1)$, we have that indeed $m_p\geq\delta|P|/q$. Further note that by the choice of each $Q_p^i$, the sets $Q_p^1,\ldots,Q_p^{m_p}$ are disjoint, as required.
\end{proof}
\subsection*{4.1 Dual SLR and their induced SLR}
Our construction of SLR sets will be via constructing and analyzing sets which we call dual SLR sets. The SLR will then be induced from these dual SLR. We start by setting some notation. Let $P$ be a non-empty finite set and $\mathbb F$ a finite field. We make use of the standard notation $\mathbb F^P$ to denote the set of all functions $\{f:P\to\mathbb F\}$. Note that $\mathbb F^P$ has a natural $\mathbb F$-vector space structure where addition is point-wise, namely, for every $f,g\in\mathbb F^P$ and $a,b\in\mathbb F$ we have that $af+bg\in\mathbb F^P$ is defined by $(af+bg)(p)=af(p)+bg(p)$ for all $p\in P$. We consider the natural inner product map $\langle\cdot,\cdot\rangle:\mathbb F^P\times\mathbb F^P\to\mathbb F$ that is defined, for $f,g\in\mathbb F^P$, by $\langle f,g\rangle=\sum_{p\in P}f(p)g(p)$. Given $f\in\mathbb F^P$, we let $f^\perp=\{g\in\mathbb F^P\mid\langle f,g\rangle=0\}$. Note that $f^\perp$ is a linear subspace of $\mathbb F^P$. More generally, given a set $S\subseteq\mathbb F^P$ we define the linear subspace $S^\perp=\bigcap_{f\in S}f^\perp$. For $f\in\mathbb F^P$ we denote $|f|=|P\setminus f^{-1}(0)|$.
For the sake of readability, the field $\mathbb F$ and the set $P$ will be omitted from the notation that we are about the introduce in this section. Both will be clear from context. For $p\in P$ define $F_p=\{f\in\mathbb F^P\mid f(p)\ne0\}$. Informally, a dual SLR is a collection of distributions over $\mathbb F^P$, one for each point $p\in P$. The distribution $D_p$, that corresponds to $p$, outputs a function $g\in\mathbb F^P$. We think of $g$ as ``passing through'' $p$. We also allow $D_p$ to output a special ``failed symbol'' $\bot$ with some small probability. A dual SLR has the guarantee that $g$ does not pass through many other points, namely, $|g|$ is bounded, and that the dimension of all functions that can be sampled, when considering all distributions $D_p$, $p\in P$, is also bounded. Perhaps most importantly is the requirement that for every other fixed $r\in P$, the sampled $g\sim D_p$ is likely to have the property that $g\notin F_r$. Formally,
\begin{ntdefinition}[Author Definition25: Dual SLR]
Let $P$ be a set, $\mathbb F$ a field. Let $D=\{D_p\mid p\in P\}$ be a collection of distributions, where for each $p\in P$, $\operatorname{supp}(D_p)\subseteq F_p\cup\{\bot\}$. Denote $S=\bigcup_{p\in P}\operatorname{supp}(D_p)$. Let $L$ be a linear subspace of $\mathbb F^P$ such that $S\subseteq L\cup\{\bot\}$. The pair $(D,L)$ is said to be a $(q,\tau,\varepsilon,\rho)$-dual SLR on $\mathbb F^P$ provided the following holds:
\begin{enumerate}
\item $|g|\leq q$ for all $g\in S\setminus\{\bot\}$.
\item For every pair of distinct $p,r\in P$ (not necessarily distinct), it holds that $\Pr_{g\sim D_p}[g(r)\ne0\mid g\ne\bot]\leq\tau$.
\item For every $p\in P$, $\Pr[D_p=\bot]\leq\varepsilon$.
\item $\dim(L)\leq(1-\rho)|P|$.
\end{enumerate}
The linear subspace $L^\perp$ of $\mathbb F^P$ is called the induced SLR from the dual SLR $(D,L)$. The parameter $\tau$ of a dual-SLR is referred to as its smoothness.
\end{ntdefinition}
Let $(D,L)$ be a dual SLR. We turn to show that, as the name suggests, the induced SLR $L^\perp$ is indeed an SLR.
\begin{ntlemma}[Author Lemma26]
Let $P$ be a set, $\mathbb F$ a field, and let $(D,L)$ be $(q,\tau,\varepsilon,\rho)$-dual SLR on $\mathbb F^P$. Then the induced SLR $L^\perp$ is a $(q-1,\tau,\varepsilon)$-SLR. Furthermore, $L^\perp$ is linear and has rate $\rho$ or larger.
\end{ntlemma}
\begin{proof}
The moreover part readily follows since $L^\perp$ is a linear subspace of $\mathbb F^P$ and since $\dim(L^\perp)=|P|-\dim(L)\geq\rho|P|$. We describe a recovering procedure for $L^\perp$, namely, a randomized algorithm that is given an oracle access to $f\in L^\perp$ as well as a point $p\in P$ as input. The recovering procedure proceeds as follows:
\begin{enumerate}
\item Sample $g\sim D_p$. If $g=\bot$ return $\bot$; Otherwise,
\item Query $f$ on all points $Q=\{r\in P\setminus\{p\}\mid g(r)\ne0\}$.
\item Return $-\frac1{g(p)}\sum_{r\in Q}g(r)f(r)$.
\end{enumerate}
The query complexity of $\operatorname{Rec}$ is bounded by $q-1$ as $|Q|=|g|-1\leq q-1$. The probability that $\bot$ is returned is at most $\varepsilon$ by construction. We turn to prove that $\operatorname{Rec}^f(p)\in\{f(p),\bot\}$. By construction, $\operatorname{Rec}^f(p)=\bot$ if and only if $g=\bot$. Assume than that $g\ne\bot$, hence, $g\in\operatorname{supp}(D_p)\subseteq L$. As $f\in L^\perp$ we have that $0=\langle f,g\rangle$, and so
\[0=\sum_{r\in P}g(r)f(r)=g(p)f(p)+\sum_{r\in Q}g(r)f(r).\]
As $g\in\operatorname{supp}(D_p)\subseteq F_p$ we have $g(p)\ne0$, and so
\[f(p)=-\frac1{g(p)}\sum_{r\in Q}g(r)f(r)=\operatorname{Rec}^f(p).\]
To conclude the proof, we turn to analyze the smoothness of $\operatorname{Rec}$. First, note that, by construction, $f$ is never queried on $p$ itself. Consider then any $r\in P\setminus\{p\}$. Conditioned on $g\ne\bot$, the function $f$ is queried on $r$ if and only if $g(r)\ne0$. Thus,
\[\Pr[f(r)\text{ is queried}]=\Pr_{g\sim D_p}[g(r)\ne0\mid g\ne\bot]\leq\tau,\]
and the proof follows.
\end{proof}
We now show that the opposite holds as well, that any linear, non-adaptive, SLR induces a dual-SLR.
\begin{ntlemma}[Author Lemma27]
Let $P$ be a set, $\mathbb F$ a field, and let $C\subseteq\mathbb F^P$ be a linear non-adaptive $(q,\tau,\varepsilon)$-SLR with rate $\rho$. Then for some set $D$, $(D,C^\perp)$ is a $(q+1,\tau,\varepsilon,\rho)$-dual SLR.
\end{ntlemma}
\begin{proof}
Let $\operatorname{Rec}$ be a recovering procedure promised by $C$ being a SLR. Assume that $\operatorname{Rec}$ uses $R$ random bits. For every point $p\in P$, denote by $E_p$ the set of choices of the random bits $r\in\{0,1\}^R$ for which $\operatorname{Rec}^c(p)=\bot$ (for any $c\in C$). Note that $\mu(E_p\subseteq\{0,1\}^R)\leq\varepsilon$.
For any $p\in P$ and $r\in\{0,1\}^R$, $r\notin E_p$, denote by $Q_{p,r}\subseteq P$ the set of query locations which $\operatorname{Rec}(p)$ makes when $r$ is the choice of randomness. Define a function $f_{p,r}:C\to\mathbb F$ such that $f(c)$ is the output of $\operatorname{Rec}^c(p)$ when fixing its randomness to $r$ and note that $f_{p,r}(c)$ only depends on $\{c(w)\mid w\in Q_{p,r}\}$, and that $f_{p,r}(c)=c(p)$. Since $C$ is linear, one can easily verify that $f_{p,r}$ is a linear map. Therefore, for some $u_{p,r}\in\mathbb F^P$, $f_{p,r}(c)=\langle u_{p,r},c\rangle$ for every $c\in C$, where $u_{p,r}(w)=0$ if $w\notin Q_{p,r}$. We have that $c(p)=\langle u_{p,r},c\rangle$ for every $c\in C$. If we define a function $g_{p,r}\in\mathbb F^P$ such that
\[g_{p,r}(w)=\begin{cases}-1,&w=p;\\u_{p,r}(q),&w\ne p,\end{cases}\]
it follows that $g_{p,r}\in C^\perp$. Note that $|g_{\{p,r\}}|\leq q+1$.
For every $p\in P$, define $D_p$ to be the following distribution. To sample from $D_p$, draw $r\in\{0,1\}^R$ uniformly at random. If $r\in E_p$, output $\bot$; otherwise, output $g_{p,r}$. Set $D=\{D_p\mid p\in P\}$ and $L=C^\perp$. It follows trivially by the definitions that $(D,L)$ is a $(q+1,\tau,\varepsilon,\rho)$-dual SLR.
\end{proof}
